% !TeX root = ../papers/nc-p3.tex
\section{三つの次数剰余類からの全次数の正則性}
\label{sec:rank-fifteen-three-residue-classes}
\begin{lemma}\label{lem:rank-fifteen-three-residue-classes}
標数零の四頂点折り畳み環 $A$ が、通常な階数15窓と同じ完全対合を持つ
標準二次表示と左右の $14641$ 複体を持つとする。
各始点の Hilbert 関数に $h_n=\binom{n+3}{3}$ の下界があり、
四中心四次式が定める $S=k[u_1,u_2,u_3,u_4]$ 上で $A$ は有限生成とする。
$M_r=\bigoplus_{n\equiv r\ (4)} A_n$ と置く。
$r=0,2,3$ で $M_r$ が $S$ 上自由で各始点の Hilbert 関数が $h_n$ なら、
残りの $M_1$ も自由で、全始点で全次数の Hilbert 関数が $h_n$ に等しい。
さらに $A$ は標準Koszul左右AS正則である。
\end{lemma}
\begin{proof}
全頂点を直和した右Koszul列の次数 $n\equiv2\pmod4$ を考える。
表示から $H_0=H_1=0$ である。上の二微分は完全対合で左の表示の双対になる。
左増大イデアル $J=A_{\ge1}$ から $A$ への次数 $d=n-4$ の左線形写像 $f$ を取る。
$u=u_1,v=u_2$ として $u f(v)=v f(u)$、$f(u),f(v)\in M_2$ を得る。
$M_2$ の自由性より $f(u)=ub$ と書ける。各矢 $a$ では $f(a)-ab\in M_3$ に対し
$u(f(a)-ab)=0$ なので、$M_3$ の自由性より $f(a)=ab$ である。
矢は左イデアルとして $J$ を生成するので $f$ は延長する。従って
$H_3=0$ であり、$H_4=0$ も $M_2$ 上で $u$ が単射であることから従う。
この列のEuler数は、各項のほかの次数が $2,0,3,2$ に属するので
\[
        \dim H_2=-4\bigl(\dim A_{n-1}-4h_{n-1}\bigr)\le0 .
\]
従って $H_2=0$、全頂点で下界が等号であり、$n-1\equiv1$ の全次数を得る。

次数 $2\bmod4$ の完全列を $S$ 加群として集める。次数シフトを省略すれば
\[
 0\longrightarrow F_4\longrightarrow F_3\longrightarrow F_2
 \longrightarrow M_1^{\oplus4}\longrightarrow F_0\longrightarrow0
\]
で、全ての $F_i$ は自由である。最後の全射は分裂するので $M_1$ の
射影次元は高々2である。四変数多項式環の斉次極大イデアルで深さは少なくとも2となる。
よって $A$ に同時に正則な二つの正次数の中心元 $u,v\in S$ を選べる。
上の延長論法で今度は $v$ が $A/uA$ に単射、$u$ が $A$ に単射であることを使うと、
全次数で $H_3=H_4=0$ を得る。左右を替えても同様である。
確定したHilbert関数のEuler等式から全次数で $H_2=0$ となり、
テンソルの双対性により左右Koszul・AS性が従う。

最後に次数 $1\bmod4$ の完全列を集めると、次数シフトを除き
$0\to M_1\to F'_3\to F'_2\to F'_1\to M_1\to0$ で中の三項は自由である。
射影次元が高々2の加群の第三シジジーは射影的なので $M_1$ は射影的であり、
次数付き自由基底を持つ。始点ごとの直和因子も自由である。
\end{proof}
\begin{corollary}
補題では不足する剰余類を任意の $r\in\mathbb Z/4$ に取っても同じ結論が成り立つ。
\end{corollary}
\begin{proof}
$n\equiv r+1$ では上の延長論法に $M_{r+1},M_{r+2}$ を使え、
Euler数で唯一未知の項は $-4\dim A_{n-1}$ である。これで同じHilbert関数を得る。
$r\ne3$ ならこの剰余類に始点単純は現れず、補題と同様に $M_r$ の射影次元は高々2である。
$r=3$ の場合は次数 $1\bmod4$ の列を使う。
$M_1,M_2$ の自由性より $H_3=H_4=0$、Euler等式から $H_2=0$ である。
未知項は中央の $M_3^{\oplus6}$ だけで、末端の自由加群への全射が分裂し、
その核への全射も分裂するので $M_3$ の射影次元は高々1である。
従って深さによる全複体の完全性が同様に従う。
$r\ne0$ では $M_r$ は自身の第三シジジーである。
$r=0$ では $K=\ker(M_0\to A_0)$ の第三シジジーとなる。
$A_0$ は $S$ 上の射影次元4であり $M_0$ の射影次元は高々2なので
$K$ の射影次元は高々3である。いずれも $M_r$ は射影的である。
\end{proof}
\begin{theorem}\label{thm:rank-fifteen-full-trace-disk}
標数零、剰余標数3の通常な階数15窓で $0<\nu(y)<3\nu(x)$ とする。
全始点で四つの中心四次式上自由階数256であり、標準Koszul左右AS正則である。
\end{theorem}
\begin{proof}
第\ref{sec:rank-fifteen-variable-zero-disk}節の $R$ と実際の四中心式で計算する。
\path{computations/verify_rank_fifteen_overlap_bound.py} は $\mathbb Q(z)$ 上の
十五桁の $H$ 級数から未切断経路と四中心式に四桁を使う。
各始点で四中心商の長さ1から7までの最大損失は $(0,0,0,2,4,2,2)$ で、
七次で文字別 $(10,10,10,10)$ の生成列と一桁の右積表が残る。
$H$ を落とす際には全係数の有理関数としての零を検査する。

右積表を $(3,z^2+z+2,H)$ で還元し、自由列で選んだ語の先頭に文字を付け、
その語を左から右へ評価して左積表も作る。第\ref{sec:rank-fifteen-overlap-annulus}節の結合則の行を使い、
長さ8から13で得る剰余体上の余次元は全始点で
\[
\begin{array}{c|cccccc}
 n&8&9&10&11&12&13\\ \hline
 \dim&31&24&10&4&1&0 .
\end{array}
\]
各段で剰余独立な行だけを持ち上げれば、その主元小行列式は単元である。
実際の商へ全射する自由な生成系を二通りの積の直和から取り、
持ち上げた独立行だけで割ると、得る両積表は計算した剰余表を持つ。
よって余次元はすべて所定の窓での上界である。
次数13の消滅と一次生成性から $A$ は $S$ 上有限である。
次数 $0,2,3\bmod4$ ではこれらの上界が $(1+t+t^2+t^3)^4$ の全係数に一致する。
各剰余類で列を持ち上げた自由 $S$ 加群の上界はちょうど $(1-t)^{-4}$ の
その剰余類部分で、各始点の一般点からの下界と等しい。従って三つの自由性が得られる。
補題\ref{lem:rank-fifteen-three-residue-classes}から全次数の結論が従う。
この議論により、実際の中心商の次数9の次元も各始点で20に決まる。
\end{proof}
\begin{corollary}\label{cor:rank-fifteen-open-integral-slopes}
ある $C(T)\in\mathbb Z[T]$ で $3\nmid C(0)$ となるものが存在する。
剰余標数が零または奇素数で、$\nu(y)>0$, $D=x^3/y$ が整、
$C(\bar D)\ne0$ なら上の全次数の結論が成立する。特に剰余標数3で
残り得る整な傾きの剰余は有限個である。
\end{corollary}
\begin{proof}
有理関数体上の全 $H$ 除算で捨てた係数は恒等的に零である。
七次までの有限個の逆元と、その後の剰余独立行を持ち上げたときの
主元の定数項は、分母と分子を合わせたある $F(z)\in\mathbb Z[z]$ を
逆にすれば全て単元である。ここに $2$ と平方根の初期分母も入れる。
$F$ は $(3,z^2+z+2)$ の外に取れる。従って $F(\bar z)\ne0$ の表示では
同じ行と列が三つの剰余類の自由性を与える。
第\ref{sec:rank-fifteen-integral-slope}節の $P,Q$ に対し、
$F(z)^{24}$ をモニックな $P(z)-TQ(z)$ で割った余りは
$\mathbb Z[T]$ 係数である。$T=0$ として3で還元すると、
$P(1+i)=0$, $F(1+i)\ne0$ なので余りは零でない。
定数項が3の単元となる余りの係数を一つ取り $C(T)$ とする。
$C(\bar D)\ne0$ なら次数24の $P-\bar DQ$ は $F^{24}$ を割らず、
従って少なくとも一つの幾何根で $F$ が消えない。その根の持上げと
$H$ の選択は同節による。
\end{proof}
\begin{corollary}\label{cor:rank-fifteen-positive-slope}
$D\ne0$ が3の上のある付値で正付値を持つとする。$x^3=Dy$ の全窓は
標準Koszul左右AS正則で、各始点で中心四次式上自由階数256である。
$D\ne27$ なら添字保持同型類は
$q^2=(1-s^2)^2-4s^6/D$ の種数2曲線から $s=0,\infty$ 上の四点を除いたもので、
全点の最小周期は2である。代数的な $D\ne0$ では、傾きの仮定は原始整数最小多項式の定数項が3の倍数であることに等しい。
\end{corollary}
\begin{proof}
付値を各点の座標体に延長する。$\nu(y)>0$ は本定理で尽くされ、
$\nu(y)=0$ は $\nu(x)>0$ なので定理\ref{thm:rank-fifteen-residue-line}、
$\nu(y)<0$ は $\nu(y)<2\nu(x)$ なので定理\ref{thm:rank-fifteen-pole-cone}を使う。
通常表示外の四点には定理\ref{thm:rank-fifteen-formal-boundary}が適用できる。
$D\ne27$ の分類と周期は第\ref{sec:rank-fifteen-gaussian-pencil}節の逆写像による。
$D=27$ は定理\ref{thm:rank-fifteen-cubic-zero-disk}の節点切片である。
最後の同値は、全ての3進素点で $\nu(D)\le0$ となることが $1/D$ の3整性に等しいことによる。
原始最小多項式の逆数多項式をモニック化した係数が全て3整である条件は、元の定数項が3の単元であることである。
\end{proof}
